% généré par scripts/06_textes.py
\dosmeta{Année de classement : 397 (directe : nom consulaire lu). Références : ICUR I 941 = ILCV +1002 = ILCV 3444 = Mossong-2022, 00030. Rattachement selon l’EDB : \emph{Incertae originis} (le volume I des ICVR réunit les inscriptions dont la provenance est perdue). Conservation selon l’EDB : \emph{s. Prassede}. Support : Plaque de marbre. 44 × 175 cm (h. × l.). Inscription gravée. Signes décrits par l’EDB : chrisme et croix monogrammatique. Reproductions signalées par l’EDB : reproduction, Perret 1851 (droits réservés ; non reproduites). Fiche EDB : Filippo Piazzolla, première version 2014. Notice du catalogue p.~\pageref{not:400}.}
\txtbloc{Texte transmis (EDCS, état 2022)}{\lnum{1}Dulcis et innoce(n)s hic dormit Severianus in somn\{i\}o pacis\par
\lnum{2}qui vixit ann⸢o⸣s p(lus) m(inus) L cuius spiritus in luce domini suscep\par
\lnum{3}tus est die VIIIIX(!) K(a)l(endas) Mart(ias) dd(ominis) nn(ostris) Flavio Caesario et Nonio Attico\par
\lnum{4}vv(iris) cc(larissimis) (con)ss(ulibus) quem locum emit Viricunda ux\{x\}or Pascasio ep(is)c(opo)}
\txtbloc{Texte de l’EDB (EDB29526)}{\lnum{1}dulcis et innoces hic dormit Severianus ((monogramma Christi)) in somnio pacis\par
\lnum{2}qui vixit annus p(lus) m(inus) L cuius spiritus in luce domini suscep/tus est ((monogramma Christi in crucem translatum)) die VIIIX k(a)l(endas) mart(ias) dd(ominis) nn(ostris) Flavio Caesario et Nonio Attico\par
\lnum{3}vv(iris) c(larissimis) c(on)s(ulibu)s quem locum emit Viricunda ux\{x\}or Pascasio EPC\par
\lnum{4}((avis cum olea in rostro)) ((monogramma Christi in crucem translatum in corona))}
